\section{La contradicción entre un hardware fijo y un Parse Tree dinámico}

Que en una imagen «el ojo pertenezca a la cara y la rueda pertenezca al coche» no es un árbol fijo; con otra imagen cambian el número de nodos, los tipos de partes y la forma en que se conectan. Un programa puede solicitar memoria de forma dinámica y escribir un pointer, pero las neuronas reales no pueden reescribir todos sus pesos de conexión en unos cuantos milisegundos. Por ello, el problema central de GLOM puede plantearse así: \emph{sin asignar dinámicamente neuron identity, ¿cómo puede la activity representar un parse tree que cambia con cada entrada?}

\lecturefigure{slide-04-hard-part-whole-hierarchies.jpg}{Cada imagen tiene un parse tree distinto, pero una red neuronal real no puede asignar neuronas nuevas al instante.}{Video oficial de Stanford Online 00:02:38--00:04:13.}

\begin{importantbox}{Contradicción central}
Los parámetros de conexión de la red $\theta$ cambian muy lentamente de un ejemplo a otro, mientras que la estructura $\mathcal T(I)$ que corresponde a la entrada $I$ es distinta cada vez. El objetivo no es que $\theta$ almacene directamente un árbol fijo, sino que el estado dinámico $\mathbf E(I)$ calculado a partir de $\theta$ codifique $\mathcal T(I)$ dentro de una sola inference.
\end{importantbox}

\subsection{Symbolic allocation, Capsule routing y universal capsule}

La respuesta de la Symbolic AI es directa: tomar un bloque de memory como node y escribir en él pointers que apuntan a otros nodes. La línea de capsules temprana, en cambio, asigna de antemano muchos nodos posibles, y cada capsule se dedica a representar un tipo de entity; una entrada solo activa unas pocas capsules, y el routing decide la parent-child connection. Esto evita crear neuronas de forma dinámica, pero trae consigo una gran cantidad de hardware inactivo, capsules type-specific y una routing optimization.

\lecturefigure{slide-05-symbolic-ai-and-capsules.jpg}{Dos líneas anteriores: memoria simbólica dinámica, o capsules asignadas de antemano seguidas de routing.}{Video oficial de Stanford Online 00:04:14--00:05:35.}

\teachervoice{Hinton valora las ideas de investigación con una frase muy útil para la enseñanza: algunas ideas «want to work» y otras no quieren funcionar; las capsules están entre ambas: con mucho esfuerzo pueden funcionar, pero el proceso es difícil. No es un rechazo de las capsules, sino un recordatorio: cuanto más fuerte es el prior estructural, más claramente debe explicarse el optimization contract.}

GLOM transforma la capsule en una \term{universal capsule}: el hardware sigue organizado en columns según la posición espacial, pero el embedding de un mismo nivel en cada column ya no queda ligado de antemano a «solo puede representar una nariz» o «solo puede representar una rueda». Puede representar cualquier identity y pose; «qué es esta vez» lo determina la activity, no el tipo fijo de la capsule.

\subsection{Column, level e island of agreement}

Se divide la imagen en posiciones espaciales $x\in\mathcal X$ y en cada posición se construye una column; dentro de la column hay varios niveles $\ell=0,1,\ldots,L$. Una misma posición puede representar, de abajo hacia arriba, textura, nostril, nose, face, person y scene. GLOM no asigna una neurona dedicada a un «face node», sino que hace que las posiciones que pertenecen a una misma face formen gradualmente, en el face level, embeddings aproximadamente idénticos.

\lecturefigure{slide-06-islands-of-identical-vectors.jpg}{Nueva representación: las islands de embeddings del mismo nivel representan los parse-tree nodes.}{Video oficial de Stanford Online 00:05:36--00:08:29.}

Sea $\mathbf e_{x,\ell}^{(t)}\in\mathbb R^d$ el vector dinámico de la posición $x$, el nivel $\ell$ y el settling step $t$. Una island candidata $S$ puede describirse con la siguiente condición aproximada:

\[
\max_{x,y\in S}\left\|\mathbf e_{x,\ell}^{(t)}-\mathbf e_{y,\ell}^{(t)}\right\|_2\le \varepsilon.
\]

Aquí, $S$ es un conjunto de posiciones en un nivel dado, $d$ es la dimensión del embedding, $t$ es el paso de iteración de la recurrent inference y $\varepsilon$ es la tolerancia de «suficientemente parecido». Las flechas bidimensionales de la clase solo usan la dirección para representar la equality; no significa que la representación real tenga solo dos dimensiones, ni exige que todas las componentes sean estrictamente bitwise identical.

\begin{importantbox}{La activity dinámica es el pointer}
Los pesos de los parámetros determinan «cómo se actualizan los vectores»; la activity del embedding determina «a quién pertenece esta posición en la entrada actual». Como la activity puede cambiar con cada imagen, una red fija puede formar islands distintas; estas islands y sus correspondencias entre niveles cumplen el papel de dynamic parse-tree nodes.
\end{importantbox}

\lecturefigure{slide-07-column-embedding-hierarchy.jpg}{Esquema de posiciones unidimensionales: en los niveles bajos los vectores son distintos; en los altos se forman gradualmente agreement islands más grandes.}{Video oficial de Stanford Online 00:08:30--00:09:15.}

\begin{knowledgebox}{Lectura de la figura: la dirección de la flecha solo codifica «si son iguales o no»}
Primero se recorren las posiciones espaciales sobre el eje horizontal y luego los niveles de abstracción sobre el eje vertical. En el nivel más bajo, las flechas negras son todas distintas, lo que indica que el contenido local de cada patch es diferente; al subir, aparecen en posiciones contiguas flechas rojas, verdes y azules con la misma dirección, lo que indica que en ese nivel se interpretan como una misma part, object o scene. Los signos de interrogación indican estados que aún no han hecho settle. La figura no proporciona los valores reales de los embeddings, su confianza ni los umbrales de frontera; por ello solo permite leer la intención del diseño y no debe leerse como un segmentation result.
\end{knowledgebox}

\begin{knowledgebox}{Términos clave: no confundir cuatro objetos en una sola palabra}
\begin{tabularx}{\textwidth}{L{0.21\textwidth}L{0.29\textwidth}X}
\toprule
Término & Qué es en GLOM & Posibles confusiones \\
\midrule
column & El hardware y el estado de varios niveles en una posición espacial & No es un channel de una CNN ni una column de una base de datos \\
level & Una escala de abstracción de la part-whole hierarchy & No es sinónimo de la capa $k$ de una red profunda común \\
embedding & El activity vector dinámico de una posición, un nivel y un paso de tiempo & No es un vector de palabras fijo ni son los pesos de la red \\
island & Región aproximadamente uniforme formada por posiciones contiguas o relacionables del mismo nivel & No es una mask dada de antemano en la entrada, y no se garantiza que coincida con el ground-truth object \\
\bottomrule
\end{tabularx}
\end{knowledgebox}

\subsection{Una estructura más amplia que la segmentación espacial ordinaria}

Hinton señala en particular que una island no tiene por qué expresar únicamente regiones contiguas de la imagen. En el lenguaje, dos fragments de una frase pueden obtener la misma representación en un nivel superior, aunque en la secuencia no formen un bloque bidimensional contiguo. Esto significa que lo que GLOM busca expresar es «pertenecer en común a una higher-level entity», y no solo una connected-component segmentation.

\teachervoice{En clase se usa el ejemplo de «shut» y «up»: en el nivel bajo son dos fragments distintos, y en el nivel alto pueden compartir el vector de «shut up». Este ejemplo sugiere que la topología de una island la determinan la neighborhood y el attention graph, y no debe entenderse en sentido estrecho como una mancha de color sólida sobre el plano de píxeles.}

\begin{warningbox}{Una island no equivale automáticamente a un object}
Aunque en algún nivel aparezcan bloques grandes de vectores parecidos, pueden provenir de la texture, del fondo, de un collapse erróneo o de un optimization shortcut. Para afirmar que «el modelo descubre objetos» se necesitan además pruebas como segmentation alignment, intervention, cross-view consistency y uso causal en tareas posteriores. Esta sesión no presenta esos experimentos.
\end{warningbox}

\subsection{Resumen de la sección}

GLOM sustituye la type-specific node allocation por universal capsules y dynamic embeddings. Su afirmación central no es «agregar más capas de attention», sino hacer que el agreement pattern de cada nivel sea en sí mismo la parse structure; el siguiente paso, el verdaderamente difícil, es lograr que ese agreement pueda formarse sin colapsar globalmente.
